Throughout, graphs are finite, simple, undirected.  The \emph{girth} $\operatorname{girth}(G)$ is the length of the shortest cycle in $G$, with $\operatorname{girth}(G)=\infty$ if $G$ is acyclic.  The chromatic number is $\chi(G)$.
We take $\log$ to be the natural logarithm.

Fix an integer $k\ge 4$.  For each $n\in\mathbb{N}$ define
\[
g_k(n)\ :=\ \max\Bigl\{\, m\in\mathbb{N}\ :\ \exists\text{ graph }G\text{ with }|V(G)|=n,\ \chi(G)=k,\ \operatorname{girth}(G)>m\,\Bigr\}.
\]
Equivalently, $g_k(n)=\max\{\operatorname{girth}(G)-1:\ |V(G)|=n,\ \chi(G)=k\}$, since $k\ge4$ forces $\operatorname{girth}(G)<\infty$.

\medskip
\noindent\textbf{Question A.} Does the limit
\[
\lim_{n\to\infty}\frac{g_k(n)}{\log n}
\]
exist (as a finite real number), and if so what is its value?

\medskip
Fix an integer $m\ge 3$.  For each $n\in\mathbb{N}$ define
\[
h^{(m)}(n)\ :=\ \max\Bigl\{\,\chi(G)\ :\ \exists\text{ graph }G\text{ with }|V(G)|=n,\ \operatorname{girth}(G)>m\,\Bigr\}.
\]

\medskip
\noindent\textbf{Question B.} Does the limit
\[
\lim_{n\to\infty}\frac{\log h^{(m)}(n)}{\log n}
\]
exist, and if so what is its value?
